\chapter{Security}\label{ch:pl:security}

On 21 February 2025 a large crypto exchange signed a transaction from one of its wallets that required several signatures. The signers
approved what their screens showed them: a legitimate transaction. The web application that prepared it had been altered two days earlier to
display legitimate transaction data while changing what was actually signed, and the signatures gave the attackers the wallet's ether. The
FBI attributed the theft of about \$1.5 billion to North Korea. No cryptography was broken; the attackers changed what people were shown.
Security in a trading firm is mostly about who and what is allowed to do things, and about making sure that what they see is what they
sign. This chapter writes the firm's access policy as data and checks it, keeps secrets and venue keys under rules, segments the network,
and looks for the insider who copies a little too much every day.

\section{What a trading firm protects, and from whom}

\begin{definition}[Threat model]\label{def:pl:security:threat}
A \emph{threat model}\index{threat model} lists what a system must protect (its assets), from whom (the adversaries and their means), through
which paths (the attack surface), and which controls address each path, so that security effort goes where the losses would be.
\end{definition}

A trading firm's assets are unusual in two ways. Some are money that moves by itself on a signature --- exchange balances, wallets (Book~3,
chapter~24), payment instructions --- so a stolen credential is a loss within minutes. Others are information whose value is its secrecy ---
strategy code, signals, parameters, positions --- so a copy is a loss even when nothing is deleted. The adversaries follow: criminals and
states for the first, competitors and departing employees for the second. The operational risk of Book~6, chapter~28, covers both; this
chapter builds the controls.

\section{Identity, access and least privilege}

\begin{definition}[Least privilege, role-based access control]\label{def:pl:security:rbac}
\emph{Least privilege}\index{least privilege} gives each person and program only the permissions its task needs, for only as long as it needs
them. \emph{Role-based access control}\index{role-based access control} grants permissions --- an action on a resource --- to roles and roles
to people, so that access follows the job rather than the individual.
\end{definition}

\begin{definition}[Multi-factor authentication]\label{def:pl:security:mfa}
\emph{Multi-factor authentication}\index{multi-factor authentication} requires more than one distinct kind of authentication factor ---
something one knows, has or is --- before access; NIST's current guidance requires, at its middle assurance level, two distinct factors and
the offer of a phishing-resistant one.
\end{definition}

The chapter's firm has twelve roles and a hundred people. Its policy is data (\texttt{firm.accessctl}): roles grant (resource, action)
pairs, people hold roles, deny rules override grants, and five segregation-of-duties constraints (Book~6, chapter~28) name pairs of
permissions no single person may hold.

\omcode{code/platforms/29-security/python/pl_accessctl.py}{22}{27}{The firm's five segregation-of-duties constraints, each a pair of
permissions that no one person may hold.}\label{lst:pl:security:sod}

\omcode{code/firm/accessctl/firm_accessctl.py}{34}{56}{The policy model: permissions through roles, minus denies; toxic combinations found at
the level of roles and of people.}\label{lst:pl:security:policy}

\begin{table}[htbp]
\centering\small
\begin{tabular}{lll}
\toprule
role & first permission & second permission\\
\midrule
developer & merge code & deploy to production\\
head of desk & propose parameters & approve parameters\\
operations & create payments & release payments\\
site reliability & create venue keys & withdraw from wallets\\
trader & submit orders & amend trades\\
\bottomrule
\end{tabular}
\caption{The five roles that, before the fix, held both halves of a segregation-of-duties constraint on their own. Through them and through
three researchers who were also risk managers, 60 of the 100 people held a toxic combination; after splitting the roles, adding a release manager and
removing the second roles, none do.}\label{tab:pl:security:toxic}
\end{table}

Each toxic role was created for convenience and each is how a known kind of loss happens: a trader who amends his own bookings can hide losses
(the rogue-trading cases of Book~6); a developer who deploys his own merge skips chapter~\ref{ch:pl:testing-and-continuous-delivery}'s review;
a head of desk who approves his own parameters skips chapter~\ref{ch:pl:signal-serving-and-parameter-management}'s four eyes. The check is a
few lines of code over the policy data and belongs in the policy's own \omterm{def:pl:testing-and-continuous-delivery:ci}{continuous integration}: a change that creates a toxic combination fails
like a failing test.

\section{Secrets and keys}

\begin{definition}[Secrets management, secret rotation]\label{def:pl:security:secrets}
\emph{Secrets management}\index{secrets management} keeps credentials --- passwords, API keys, private keys, database passwords --- in a
vault that issues them to authenticated programs under leases and logs every access, never in code or configuration files. \emph{Secret
rotation}\index{secret rotation} replaces a secret on a schedule and on any suspicion, so that a leaked secret is valid only until its next
rotation: its cryptoperiod, in NIST's terms.
\end{definition}

\begin{definition}[Hardware security module]\label{def:pl:security:hsm}
A \emph{hardware security module}\index{hardware security module} (HSM) is a tamper-resistant device that generates and stores keys and
performs operations with them, so that private keys never leave it in plain form; signing happens inside, on data the host presents.
\end{definition}

Rotation bounds exposure: a secret that leaks at a random time and is rotated every 90 days stays valid 45 days on average; rotated weekly,
three and a half. An HSM, or the multi-party computation of Book~3, chapter~24, removes the key from the host --- but not the question the hook
asks: the device signs whatever it is shown. The defence is to make what is signed verifiable independently of the screen that prepared it:
a second channel that shows the transaction's actual content, a policy engine that refuses a transaction type (a contract upgrade, a new
destination) whatever the signers approve.

\begin{omfigure}
\begin{tikzpicture}[font=\small,node distance=6mm,
  box/.style={draw=omDef,fill=omDef!8,rounded corners=2pt,align=center,minimum height=11mm,text width=31mm},
  bad/.style={draw=omThm,fill=omThm!10,rounded corners=2pt,align=center,minimum height=11mm,text width=24mm},
  arr/.style={-{Stealth[length=2mm]},thick}]
\node[bad] (app) {web application\\(altered)};
\node[box,right=18mm of app,yshift=10mm] (scr) {signer's screen:\\the expected transfer};
\node[box,right=18mm of app,yshift=-10mm] (dev) {signing device:\\another transaction};
\node[box,right=10mm of dev,text width=27mm] (chk) {independent decoder\\and policy engine};
\node[box,right=10mm of chk,text width=14mm] (chain) {wallet};
\draw[arr] (app) -- node[above,sloped,font=\scriptsize]{what is shown} (scr);
\draw[arr,omThm] (app) -- node[below,sloped,font=\scriptsize]{what is signed} (dev);
\draw[arr] (dev) -- (chk);
\draw[arr] (chk) -- node[below,font=\scriptsize,align=center]{only if\\allowed} (chain);
\draw[dashed,gray] (scr.east) -| (chk.north);
\node[gray,font=\scriptsize,anchor=south west] at (scr.north east) {the check made on the screen};
\end{tikzpicture}

\omcaption{The hook's attack in one picture: the altered application showed the signers one transaction and sent their devices another. A
control that decodes what is actually to be signed, independently of the application, and a policy that refuses what no signer should approve
(a change to the wallet's own contract, a new destination), check the thing signed rather than the thing shown.}\label{fig:pl:security:signing}
\end{omfigure}

\section{Exchange and wallet credentials}

\begin{definition}[Withdrawal allow-list]\label{def:pl:security:allowlist}
A \emph{withdrawal allow-list}\index{withdrawal allow-list} restricts the destinations to which an account or key may send assets to a list
fixed in advance, changed only through a separate, slower procedure, so that a stolen credential cannot move assets to the thief.
\end{definition}

Venue API keys (Book~3, chapter~15) are the firm's most exposed secrets: they sit on the trading servers and sign every order. The chapter's
keys carry scopes --- a trading key can read and trade but not withdraw; a withdrawal key can only withdraw, to allow-listed addresses --- and
every request is signed with Book~13's HMAC signer and stamped with its time. The venue's checks decide what a stolen key can do
(\cref{tab:pl:security:keys}).

\omcode{code/firm/accessctl/firm_accessctl.py}{99}{114}{The venue's checks on a request: a known key, a valid signature, a fresh timestamp, an
action within the key's scope, and a withdrawal address on the allow-list.}\label{lst:pl:security:verify}

\begin{table}[htbp]
\centering\small
\begin{tabular}{lll}
\toprule
attempt & accepted & reason\\
\midrule
trade with the trading key & yes & ok\\
withdraw with the trading key & no & action outside the key's scope\\
withdraw to an outside address & no & address not on the allow-list\\
withdraw to the custodian & yes & ok\\
forged signature & no & bad signature\\
replayed a minute later & no & stale request\\
\bottomrule
\end{tabular}
\caption{A stolen trading key and a stolen withdrawal key tried against the venue's checks. The thief can trade --- which can still cost money
--- but cannot move assets anywhere the firm has not listed in advance.}\label{tab:pl:security:keys}
\end{table}

\section{Segmentation and insider risk}

\begin{definition}[Network segmentation, zero-trust architecture]\label{def:pl:security:segmentation}
\emph{Network segmentation}\index{network segmentation} divides a firm's network into zones --- trading, research, corporate, the outside ---
with traffic between zones allowed only on stated paths, so that a compromise in one zone does not reach the others. A \emph{zero-trust
architecture}\index{zero-trust architecture} grants no implicit trust for being inside a network: every request is authenticated and
authorised on its own, by identity and device, before a session to a resource is established.
\end{definition}

\begin{omfigure}
\begin{tikzpicture}[font=\small,
  zone/.style={draw=omDef,fill=omDef!6,rounded corners=3pt,minimum width=30mm,text width=30mm,minimum height=16mm,align=center},
  hot/.style={draw=omThm,fill=omThm!8,rounded corners=3pt,minimum width=30mm,text width=30mm,minimum height=16mm,align=center},
  arr/.style={-{Stealth[length=2mm]},thick}]
\node[hot] (t) at (0,0) {trading zone\\\scriptsize strategies, gateways, keys};
\node[zone] (r) at (6.2,0) {research zone\\\scriptsize tick store, notebooks};
\node[zone] (c) at (12.2,0) {corporate zone\\\scriptsize e-mail, laptops};
\node[zone] (v) at (0,-2.8) {venues\\\scriptsize cross-connects};
\node[zone] (i) at (12.2,-2.8) {internet};
\draw[arr] (t) -- node[left,font=\scriptsize]{orders only} (v);
\draw[arr] ([yshift=3mm]t.east) -- node[above,font=\scriptsize]{fills, positions} ([yshift=3mm]r.west);
\draw[arr,dashed] ([yshift=-3mm]r.west) -- node[below,font=\scriptsize,align=center]{parameters, through\\review (ch.~16)} ([yshift=-3mm]t.east);
\draw[arr] (c) -- node[above,font=\scriptsize,align=center]{authenticated\\access} (r);
\draw[arr] (i) -- node[right,font=\scriptsize]{mail, web} (c);
\draw[omThm,thick] (6.2,-2.8) node[cross out,draw,minimum size=4mm]{} node[below=3mm,font=\scriptsize,align=center]{no path from the\\internet or laptops to trading};
\end{tikzpicture}

\omcaption{The chapter's network zones and the only paths allowed between them. The trading zone, which holds the venue keys, talks to venues
and sends fills to research; parameters come back only through chapter~16's reviewed path; nothing from the corporate zone or the internet
reaches it directly.}\label{fig:pl:security:zones}
\end{omfigure}

\begin{definition}[Insider threat, data-loss prevention]\label{def:pl:security:insider}
An \emph{insider threat}\index{insider threat} is harm done by someone with legitimate access --- an employee, a contractor --- through theft,
sabotage or negligence. \emph{Data-loss prevention}\index{data-loss prevention} is the set of controls that detect and block the movement of
protected data out of the places it is allowed to be: copies to removable media, uploads, mail, and unusual volumes of reading.
\end{definition}

The best-known case in trading is also a warning about the law's limits: in 2009 a programmer uploaded more than 500\,000 lines of his
employer's high-frequency trading code to a server in Germany on his last day; his federal conviction was reversed on appeal in 2012, on
the reach of the statutes rather than the facts. A firm cannot rely on the courts (Book~16, chapter~11, on trade secrets); it must see the
copying. The difficult insider is not the one who takes everything on the last day but the one who takes a little more than usual every day.

The chapter's model gives each of the hundred people a year of 250 trading days of access logs --- the strategy-code files each reads a
day, at each person's own level, with a weekly rhythm and noise --- and plants an insider who, from a random day, reads a stated share more
than usual, every day. Four detectors score each person against their own baseline, the median and median absolute deviation of 30 days:
a single-day threshold, or a CUSUM of the daily scores (Book~7, chapter~13), each with a baseline that trails today or one that ends 60 days
earlier. All four are calibrated on clean people to one false alarm a month per hundred people.

\omcode{code/firm/accessctl/firm_accessctl.py}{118}{137}{Per-person robust scores against a baseline that can be lagged, and the one-sided CUSUM
that accumulates small daily excesses.}\label{lst:pl:security:scores}

\begin{omfigure}
\begin{tikzpicture}
\begin{axis}[width=11.5cm,height=5.6cm,xlabel={extra reading per day (\% above the person's usual, log scale)},ylabel={insiders flagged (of 20)},
  grid=major,grid style={gray!20},tick label style={font=\small},label style={font=\small},table/col sep=comma,
  xmode=log,xmin=4,xmax=120,ymin=0,ymax=21,xtick={5,10,20,30,50,100},xticklabels={5,10,20,30,50,100}]
\addplot[omThm,thick,mark=*] table[x=extra_pct,y=cusum_lagged_found]{figdata/platforms/29-security/detection.csv};
\addplot[omThm,thick,dashed,mark=o] table[x=extra_pct,y=cusum_trailing_found]{figdata/platforms/29-security/detection.csv};
\addplot[omDef,thick,mark=square*] table[x=extra_pct,y=day_lagged_found]{figdata/platforms/29-security/detection.csv};
\addplot[omDef,thick,dashed,mark=square] table[x=extra_pct,y=day_trailing_found]{figdata/platforms/29-security/detection.csv};
\node[omThm,font=\scriptsize,anchor=south east] at (axis cs:28,17) {CUSUM, lagged baseline};
\node[omThm,font=\scriptsize,anchor=north west] at (axis cs:32,6) {CUSUM, trailing};
\node[omDef,font=\scriptsize,anchor=south east] at (axis cs:48,12.5) {single day, lagged};
\node[omDef,font=\scriptsize,anchor=south] at (axis cs:45,1.6) {single day, trailing};
\end{axis}
\end{tikzpicture}

\omcaption{Planted insiders flagged, out of twenty, against how much more than usual they read each day, at one false alarm a month per
hundred people. A single-day threshold on a trailing baseline almost never sees them; a CUSUM on a baseline lagged by 60 days finds 12 of 20 at
+20\% a day and 19 of 20 at +30\%. Data: \texttt{fig\_accessctl.py}.}\label{fig:pl:security:detection}
\end{omfigure}

Two things matter, and \cref{fig:pl:security:detection} separates them. A single day of reading 30\% more than usual is well inside a person's
noise; only accumulation --- the CUSUM --- turns thirty such days into a signal. And a baseline that trails today learns the new habit: after
thirty days the insider's copying is the person's normal, and the score returns to zero. Lagging the baseline by sixty days keeps the copying
out of its own reference long enough to be seen. With both, an insider reading 30\% more than usual is flagged in a median of 22 days (19 of
20), one reading 50\% more in 12 (all 20); a single-day threshold on a trailing baseline flags one of twenty.

\begin{omfigure}
\begin{tikzpicture}
\begin{axis}[width=11.5cm,height=5.2cm,xlabel={trading day},ylabel={CUSUM of daily scores},
  grid=major,grid style={gray!20},tick label style={font=\small},label style={font=\small},table/col sep=comma,
  xmin=90,xmax=250,ymin=0,ymax=60]
\addplot[omThm,thick] table[x=day,y=cusum_lagged]{figdata/platforms/29-security/insider.csv};
\addplot[omThm,dashed] table[x=day,y=cusum_trailing]{figdata/platforms/29-security/insider.csv};
\addplot[gray,thin] table[x=day,y=clean_cusum_lagged]{figdata/platforms/29-security/insider.csv};
\draw[black,dashed] (axis cs:90,12.709) -- (axis cs:250,12.709) node[pos=0.98,above left,font=\scriptsize]{threshold 12.7};
\draw[black,dotted] (axis cs:120,0) -- (axis cs:120,40) node[pos=0.75,left,font=\scriptsize,align=right]{copying\\starts};
\node[omThm,font=\scriptsize,anchor=west] at (axis cs:205,57) {lagged baseline};
\node[omThm,font=\scriptsize,anchor=west] at (axis cs:175,6) {trailing baseline};
\end{axis}
\end{tikzpicture}

\omcaption{One planted insider reading 30\% more strategy code than usual from day 120. The CUSUM on the lagged baseline crosses its calibrated
threshold on day 131; on the trailing baseline it just reaches its own threshold (11.2) on day 147, then falls back once the copying has entered
the baseline. A clean person's
CUSUM (grey) stays near zero. Data: \texttt{fig\_accessctl.py}.}\label{fig:pl:security:insider}
\end{omfigure}

\begin{dated}{2026-09}{The cases and the standards}\label{dat:pl:security:sources}
The FBI's public service announcement of 26 February 2025 attributed to North Korea the theft of about \$1.5 billion in virtual assets from the
exchange Bybit on or about 21 February; the investigation published by Sygnia describes JavaScript altered on 19 February in the Safe\{Wallet\}
web application so that its interface displayed legitimate transaction data while the transaction itself was altered. In \emph{United States
v.\ Aleynikov} (2d Cir.\ 2012) the court reversed a conviction for uploading more than 500\,000 lines of trading source code. NIST SP 800-207
(2020) defines \omterm{def:pl:security:segmentation}{zero-trust architecture}; SP 800-63B (revision of August 2025) the authentication assurance levels; SP 800-57 Part 1 (2020) the
cryptoperiod of a key.
\end{dated}

\section{Tutorial: the slow copy}\label{tut:pl:security:tut}

\begin{tutorial}
\textbf{Goal.} Check the access policy for toxic combinations, issue and test scoped venue keys, and find a slow insider at a fixed false-alarm
rate. \textbf{End state:} \cref{tab:pl:security:toxic,fig:pl:security:detection}.
\begin{tutsteps}
\item \textbf{Policy}: \texttt{pl\_accessctl.policies()} and \texttt{toxic\_summary()}, before and after the fix.
\item \textbf{Keys}: \texttt{stolen\_key\_attempts()}; add a key with a trade scope and an allow-list and try it.
\item \textbf{Vault}: \texttt{firm\_accessctl.Vault} with a 90-day lease; rotate and list what expires.
\item \textbf{Logs}: \texttt{access\_logs()} and \texttt{statistics(x)}; \texttt{thresholds()}.
\item \textbf{Insider}: \texttt{delays(extra)} for each share of extra reading.
\end{tutsteps}
\textbf{What to change next.} Score people against their peers in the same role as well as against themselves; add file-level features (files
never read before) and see how much sooner the insider is flagged.
\end{tutorial}

\section{Build: access control}\label{bld:pl:security:accessctl}

\begin{build}
\bfield{Purpose} Grant each person and program only what its job needs, keep secrets and venue keys under rules that bound a leak, and see a
slow copy at a stated false-alarm rate.
\bfield{Interface} \texttt{Policy} (\texttt{allowed}, \texttt{permissions}, \texttt{toxic\_roles}, \texttt{toxic\_users}), \texttt{Vault},
\texttt{VenueKey}, \texttt{request}, \texttt{verify}, \texttt{robust\_scores}, \texttt{cusum}, \texttt{onsets}, \texttt{calibrate},
\texttt{first\_alarm}.
\bfield{Rules} Policy as reviewed data, checked for toxic combinations in \omterm{def:pl:testing-and-continuous-delivery:ci}{continuous integration}; secrets only from the vault, under leases,
rotated; venue keys scoped, withdrawal only to allow-listed addresses; every request signed and time-stamped; detectors calibrated on clean
populations and reviewed monthly.
\bfield{Acceptance tests} \texttt{code/firm/accessctl/tests/}: grants, denies, toxic roles and users; vault leases, rotation and expiry; each of
the venue's checks; lagged robust scores, CUSUM, calibration to a false-alarm rate, and a persistent shift found.
\bfield{Stretch} Just-in-time access grants that expire; peer-group scoring; an independent transaction decoder for signers.
\end{build}

\begin{omsources}
\item Federal Bureau of Investigation, \emph{North Korea Responsible for \$1.5 Billion Bybit Hack}, PSA 250226, 26 February 2025; Sygnia,
investigation of the Bybit hack, 2025.
\item \emph{United States v.\ Aleynikov}, 676 F.3d 71 (2d Cir.\ 2012).
\item NIST SP 800-207, \emph{Zero Trust Architecture}, 2020; SP 800-63B, \emph{Authentication and Authenticator Management}; SP 800-57 Part 1
Rev.~5, \emph{Recommendation for Key Management}, 2020.
\item One Quant Book~3, chapters~14, 15 and~24 (keys, venue APIs, wallets); Book~6, chapter~28 (segregation of duties); Book~16, chapter~11
(trade secrets).
\end{omsources}

\section{Exercises}

\begin{exercise}[$\star$]\label{exo:pl:security:1}
Why does the check for toxic combinations look at people as well as roles?
\end{exercise}

\begin{exercise}[$\star$]\label{exo:pl:security:2}
A secret leaks at a random moment. How long does it stay valid on average under 90-day and weekly rotation?
\end{exercise}

\begin{exercise}[$\star$]\label{exo:pl:security:3}
What can a thief do with the firm's stolen trading key, and what can he not?
\end{exercise}

\begin{exercise}[$\star\star$]\label{exo:pl:security:4}
Why does a baseline that trails today stop seeing a slow copy after about a month?
\end{exercise}

\begin{exercise}[$\star\star$]\label{exo:pl:security:5}
Why is a single-day threshold the wrong tool for a slow copy, whatever its level?
\end{exercise}

\begin{exercise}[$\star\star$]\label{exo:pl:security:6}
The signers of the hook each checked the transaction on their screen. What control would have caught the attack?
\end{exercise}

\begin{exercise}[$\star\star\star$]\label{exo:pl:security:7}
\emph{Coding.} Rerun \texttt{delays(0.3)} with the CUSUM's reference value \texttt{K} at 0.25 instead of 0.5 (recalibrating the thresholds).
What happens to the detections and why?
\end{exercise}

\begin{exercise}[$\star\star\star$]\label{exo:pl:security:8}
\emph{Find the flaw.} \textquotedbl{}Our trading network is behind the firewall, so the systems inside it can trust each other.\textquotedbl{}
\end{exercise}

\section{Problem: The Slow Copy}

\begin{problem}[{Weekend problem --- the slow copy}]\label{pb:pl:security:1}
The chapter's policy, keys, vault and access logs.

\medskip\noindent\textbf{Part I --- Access.}
\begin{enumerate}
\item What does the firm protect, and from whom?
\item What are the five segregation-of-duties constraints?
\item Which roles were toxic before the fix, and how many people held a toxic combination?
\item What did the fix change?
\item Where should the toxic-combination check run?
\end{enumerate}

\medskip\noindent\textbf{Part II --- Secrets and keys.}
\begin{enumerate}[resume]
\item What does a vault add to a secret in a configuration file?
\item How does rotation bound a leak?
\item What are the venue's checks on a request?
\item What does each stolen key allow?
\item What does an HSM protect, and what does it not?
\end{enumerate}

\medskip\noindent\textbf{Part III --- Insiders.}
\begin{enumerate}[resume]
\item How are the four detectors built and calibrated?
\item What does each flag at +30\% a day?
\item Why does the CUSUM matter?
\item Why does the lag of the baseline matter?
\item When does the example insider's lagged CUSUM cross its threshold?
\end{enumerate}

\medskip\noindent\textbf{Part IV --- The verdict.}
\begin{enumerate}[resume]
\item State the \emph{named result}: the days until a planted insider's slow exfiltration is flagged at one false alarm a month per hundred
users, and the number of toxic role combinations in the firm's policy before and after the fix.
\item What would you add to the detectors first?
\item How should the trading zone be segmented?
\item What is the lesson of the hook for any firm that signs transactions?
\item In one sentence: what is security in a trading firm about?
\end{enumerate}
\end{problem}

\section{Interview questions}

\begin{interviewq}[$\star$ \iqroles{developer}]\label{iq:pl:security:1}
Where should an exchange API secret live, and how should a strategy get it?
\end{interviewq}

\begin{interviewq}[$\star\star$ \iqroles{developer}]\label{iq:pl:security:2}
What permissions should an exchange API key used for trading have?
\end{interviewq}

\begin{interviewq}[$\star\star$ \iqroles{developer}]\label{iq:pl:security:3}
How would you detect an employee who copies strategy code slowly over months?
\end{interviewq}

\begin{interviewq}[$\star\star$ \iqroles{developer}]\label{iq:pl:security:4}
What is zero trust, and what does it change inside a trading network?
\end{interviewq}

\begin{interviewq}[$\star\star\star$ \iqroles{developer}]\label{iq:pl:security:5}
Design access control, secrets and monitoring for a trading firm of a hundred people with venue keys and a crypto treasury.
\end{interviewq}
